\documentclass[12pt,a4paper]{report}

%------------------------------------------------
% Packages
%------------------------------------------------
\usepackage[T1]{fontenc}
\usepackage[utf8]{inputenc}
\usepackage{lipsum}        % For placeholder text
\usepackage{hyperref}      % Clickable references
\usepackage{graphicx}      % Images
\usepackage{amsmath}       % Math support
\usepackage{amssymb}
\usepackage{setspace}      % Line spacing
\usepackage{geometry}      % Page margins
\usepackage{graphicx}
\geometry{margin=1in}

%------------------------------------------------
% Document
%------------------------------------------------
\begin{document}

%------------------------
% Title Page
%------------------------
\begin{titlepage}
    \centering

    {\Huge \textbf{Adaptivity report}\par}
    \vspace{1.5cm}
    \date{\today}

    {\Large Adam Adaptiosson\par}
    \vspace{0.5cm}

    \vfill

    % --- Your custom text here ---
    \begin{center}
    \textbf{Weld control log:}\\[4pt]
    \texttt{\detokenize{{log_file}}}
    \end{center}

    \vspace{1cm}

    % optional footer
    {\small ESAB Automation Adaptio Team}

\end{titlepage}


\newpage

%------------------------------------------------
\section{Profile data}
The following deposition sequences are constructed from laser scanner profiles and then color coded 
based on the logged welding parameters.

\subsection{Color scheme}
Parameters are mapped linearly to color. 64 colors are used. Parameter ranges are configurable. If the parameters of a
bead are mapped to outside a valid range they will have get the color white (below valid range) or the color black (above valid range).

\subsection{Coordinate system}
The coordinate system used is the weld object coordiante system (ROCS) with origin on the rotation axis. The plot coordinate are
cylindrical, so the MCS x and z coordinates of the logged profiles have been transformed to ROCS and converted to cylindrical. 
The plots basically show radial cross-sections.

\subsection{Torch position}
The wire tip position for each bead and position is plotted as a circle with white background. The exact position should be interpreted with 
caution since it requires that the stickout configured in the plot script is identical to the stickout used during welding.

\subsection{Bead detection}
In order to reduce the occurence of near but not perfect overlapping scan profiles the scan data is filtered through
a bead detection algorithm. This algorithm tries to identify the contiguous region associated with a bead. Within the bead region
the two subsequent profiles are expected to differ significantly. Outside the bead region the profiles should ideally be identical, but in reality
they differ slightly due to inaccuracy in position etc. The bead detection algorithm creates a filtered version of the current profile. Points in
the bead region is taken from the current profile and points outside the bead region is taken from the reference (previous) profile.
It is important to note that flux grains present in the profiles appear as spikes and these can distrupt the bead detection algorithm.

\subsection{Bead switch region}
In the bead switch region, where the torch crosses over from one wall to the other, the plots will be difficult to interpret. This usually
occurs around position = 0 degrees.

\subsection{Overlap region}
After welding a full revolution, the torch is not moved immediately to a new position. There is an overlap such that weld metal will 
be deposited on top of the just welded bead. This overlap region is a few cm long and can sometimes be seen in the plots.

\subsection{Last bead}
If weld movement is stopped when the last bead is finshed, there will be no scan data for the last bead. This is because the scanner reads
the joint before the torch and therefore have not "seen" the last bead yet.

\subsection{Blind spot}
The region between the laser line and the torch at the start of welding will never be observed by the scanner. This part of the weld is missing
from the logs. This will appear as a strange joint bottom for positions close to 360 degrees. In this region, what is actually the first bead,
will be interpreted as joint bottom and therefore colored gray.

\subsection{Notes on accuracy}
During acqusition of scanner data, the weld object drifts and deforms due to heat etc. This means that the profiles from subsequent
revolutions will not perfectly align and can therefore not be plotted on top of each other in a meaningful way.
In order to be able to plot all the beads aligned w.r.t the "ideal" joint, the laser profiles are translated, rotated and deformed
slightly. Usually these profile alterations are of a sub-millimeter magnitude. 

Another source of error is the mapping of logged welding parameters to the logged laser profiles. Since the scanner observes the
joint almost a full revolution after the deposition has been made, the corresponding welding parameters must be gotten from another
location in the log. In order to do this, the log parser must have been given information about the weld object calibration so that
it can compute the rotation angle between laser line and torch plane.


{figures}

\end{document}
